\section{Admissible scenes and well-posedness assumptions}
\label{sec:admissible-scenes}

The phrases arbitrary pose, arbitrary package size, and arbitrary medium mean
that these quantities are not restricted to one fixed parameter chart or world
grid. They do not mean that physically singular or mathematically ill-posed
scenes must produce a finite prediction.

\subsection{Geometric admissibility}

A scene is admissible when:
\begin{enumerate}
\item every object map is at least piecewise $C^1$ with finite measure and a
      well-defined outward normal almost everywhere where a surface formulation
      is used;
\item conductor cross-section remains positive and the sweep Jacobian does not
      vanish;
\item distinct perfectly conducting volumes do not overlap unless an explicit
      electrical contact/junction is declared;
\item material interfaces have a declared ownership rule and do not create
      ambiguous multiply occupied volumes;
\item terminal regions have nonzero measure and define a valid oriented
      work-conjugate port pair.
\end{enumerate}

Near-touching objects are allowed, but the numerical backend must refine
near-singular quadrature accordingly. Exact topological contact is a different
scene and requires an explicit contact model.

\subsection{Material admissibility}

Over the declared frequency/temperature domain, constitutive laws are bounded,
causal, and passive. For a simple isotropic conductive dielectric under a fixed
phasor convention this implies non-negative dissipative response. For tensor
media, the Hermitian dissipative part is positive semidefinite.

Thermal density and heat capacity remain positive:
\[
 \rho_m(\mathbf r,T)>0,\qquad c_p(\mathbf r,T)>0,
\]
and the symmetric part of thermal conductivity is uniformly positive on each
conducting thermal region:
\[
 \xi^\T k(\mathbf r,T)\xi
 \ge k_{\min}\norm{\xi}^2,\qquad k_{\min}>0.
\]
Degenerate phase change, hysteresis, superconductivity, plasma breakdown, and
other constitutive phenomena require an explicit material extension rather than
being hidden inside extrapolation.

\subsection{Frequency domain}

Every artifact declares a finite operating band
\[
 \omega\in[\omega_{\min},\omega_{\max}].
\]
The REFERENCE formulation may cover DC/quasi-static, induction-dominated, and
full-wave regimes, but the basis and preconditioner policy may change between
regimes. Numerical formulations that suffer interior resonances must use a
combined/stabilized integral formulation so a physical scattering solution is
not confused with a discretization null space.

\subsection{Thermal boundary well-posedness}

A pure Neumann insulated thermal problem has an undetermined absolute
temperature under zero net heat removal and no finite steady state under
positive continuous heating. Such a scene may still have a transient solution,
but a steady-state query must return \texttt{NO\_FINITE\_STEADY\_STATE} rather
than inventing a stabilized answer.

\subsection{Nonlinear operating domain}

For temperature-dependent coupled simulation, the trajectory must remain within
the declared constitutive temperature range. Crossing the range triggers
extrapolation status, model escalation, or simulation termination according to
policy.

\subsection{Existence of certificates}

CERTIFIED mode assumes that the chosen physical operator admits a computable
residual and a usable stability/goal estimator in the requested regime. If the
operator is too ill-conditioned for a meaningful bound, the system may still
return a REFERENCE numerical result but must not label it certified.

\begin{principle}[Explicit invalidity beats hidden regularization]
Geometry repair, artificial conductivity floors, thermal leakage, or numerical
damping may be used only when they are declared physical/numerical model
choices. They are never silently introduced to force an invalid scene to
converge.
\end{principle}
